\chapter{Main theorem}\label{chap:main}

Combining the envelope domination from Chapter~\ref{chap:gaussian-comparator}
with the extremal witness from Chapter~\ref{chap:envelope} yields the sharp
comparison theorem.

\begin{theorem}[Sharp one-dimensional convex sub-Gaussian comparison]\label{thm:main}
\lean{sharp_convexSubgaussianComparison}\leanok
\uses{thm:cx-route,thm:extremizer-stoploss,thm:gauss-ge-J}
The sharp admissible Gaussian scale is $c_\star = c_0$. Equivalently:
\begin{enumerate}
\item Every centered real random variable satisfying
  $\Prb(|X|>t)\le 2e^{-t^2/2}$ for all $t \ge 0$ is convex-dominated by the
  centered Gaussian law of scale $c_0$.
\item For every $c < c_0$, convex domination by the Gaussian law of scale $c$
  fails, already for the canonical witness $\mu^\star$ and an explicit hinge
  test function.
\end{enumerate}
\end{theorem}

\begin{proof}\leanok
For domination: Theorem~\ref{thm:gauss-ge-J}
shows that the Gaussian stop-loss transform at scale $c_0$ dominates the sharp
envelope $J_G$. Theorem~\ref{thm:cx-route} then upgrades this stop-loss
comparison to full convex domination.

For sharpness: Theorem~\ref{thm:extremizer-stoploss} identifies the stop-loss of
the canonical witness $\mu^\star$ with the envelope itself. At the
threshold $u = cz$, the Gaussian stop-loss of scale $c$ evaluates to
$c\varphi(z) - up_0$, while the envelope is at least $B - up_0$. If
$c < c_0 = B/\varphi(z)$, then the witness stop-loss is strictly larger. The
hinge function $x \mapsto (x-u)_+$ therefore exhibits a convex test
function for which the comparison fails.
\end{proof}

For code navigation, the two halves also appear separately in
\texttt{GaussianAsymptotics.lean} as \texttt{gaussianDomination\_c0} and
\texttt{gaussianSharpness\_below\_c0}.
